\DocumentMetadata{
  pdfversion=2.0,
  pdfstandard=ua-2,
  lang=en-US,
  tagging=on,
  tagging-setup={math/setup=mathml-SE,tabsorder=structure}
}
\documentclass[11pt,letterpaper]{article}
\usepackage[margin=0.68in,includefoot]{geometry}
\usepackage{newcomputermodern}
\usepackage{amsmath,amscd,mathtools}
\usepackage{tikz,adjustbox}
\providecommand{\square}{\mdlgwhtsquare}
\usepackage[hidelinks]{hyperref}
\hypersetup{
  pdftitle={Logic PhD Exam, September 1989},
  pdfauthor={University of Florida Department of Mathematics},
  pdfsubject={Graduate mathematics examination},
  pdfdisplaydoctitle=true,
  bookmarksopen=true
}
\setcounter{secnumdepth}{0}
\setlength{\parindent}{0pt}
\setlength{\parskip}{0.28em}
\setlength{\emergencystretch}{3em}
\setlength{\leftmargini}{2.15em}
\setlength{\leftmarginii}{2.65em}
\setlength{\leftmarginiii}{3em}
\pagestyle{empty}
\raggedbottom
\allowdisplaybreaks
\NewTaggingSocketPlug{block/list/label}{examproblem}{%
  \tagstructbegin{tag=Lbl}\tagstructend%
  \tagstructbegin{tag=\LBody}%
  \tagstructbegin{tag=\ProblemHeadingTag,title-o={\ProblemHeadingText}}%
  \tagmcbegin{tag=\ProblemHeadingTag,actualtext={\ProblemHeadingText}}%
  #2%
  \tagmcend%
  \tagstructend%
}
\newenvironment{examproblems}[2]{%
  \def\ProblemHeadingTag{#2}%
  \AssignTaggingSocketPlug{block/list/label}{examproblem}%
  \begin{enumerate}%
  \setcounter{enumi}{#1}%
  \renewcommand{\labelenumi}{\textbf{\theenumi.}}%
  \setlength{\topsep}{0.35em}%
  \setlength{\partopsep}{0pt}%
  \setlength{\itemsep}{0.48em}%
  \setlength{\parsep}{0.18em}%
}{\end{enumerate}}
\newenvironment{examsubparts}{%
  \AssignTaggingSocketPlug{block/list/label}{default}%
  \begin{enumerate}%
  \setlength{\topsep}{0.18em}%
  \setlength{\partopsep}{0pt}%
  \setlength{\itemsep}{0.16em}%
  \setlength{\parsep}{0.1em}%
}{\end{enumerate}}
\newenvironment{examinstructions}{%
  \par\begingroup\itshape
}{%
  \par\endgroup\vspace{0.18em}
}
\makeatletter
\renewcommand\subsection{\@startsection{subsection}{2}{0pt}%
  {-1.25ex plus -.25ex minus -.1ex}%
  {0.55ex plus .1ex}%
  {\normalfont\large\bfseries}}
\renewcommand{\@maketitle}{%
  \newpage\null\vskip -1em%
  {\footnotesize\bfseries
    \href{https://www.ufl.edu/}{University of Florida}, \href{https://math.ufl.edu/}{Department of Mathematics}\par}%
  \vskip -0.5em%
  \tagstructbegin{tag=H1,title={Logic PhD Exam, September 1989}}%
  \tagpdfparaOff%
  \tagmcbegin{tag=H1}%
    {\LARGE\bfseries\href{https://gma.math.ufl.edu/exams/logic-phd/}{Logic PhD Exam}, September 1989\par}%
  \tagmcend%
  \tagpdfparaOn%
  \tagstructend%
  \vskip 0.25em%
  \tagmcbegin{artifact}%
    \hrule height 1pt%
  \tagmcend%
  \vskip 1em%
}
\makeatother
\title{Logic PhD Exam, September 1989}
\author{}
\date{}
\begin{document}
\maketitle
\thispagestyle{empty}
\begin{examinstructions}
Do 2 from each section.
\end{examinstructions}
\subsection{I. General Logic}
\begin{examproblems}{0}{H3}
\def\ProblemHeadingText{Problem 1.}
\item
State and prove the Löwenheim–Skolem Theorem for countable languages.
\def\ProblemHeadingText{Problem 2.}
\item
State and prove Herbrand’s Theorem.
\def\ProblemHeadingText{Problem 3.}
\item
Suppose that~\(\mathcal{N}\) is a submodel of~\(\mathcal{M}\). Prove that if~\(A\) is a quantifier-free sentence, then
\[
\mathcal{N} \models A \quad\Longleftrightarrow\quad \mathcal{M} \models A.
\]
\end{examproblems}
\subsection{II. Set Theory}
\begin{examproblems}{3}{H3}
\def\ProblemHeadingText{Problem 4.}
\item
Prove the Mostowski Collapsing Lemma: If \((M,E)\) is a well-founded model of extensionality, then there is a transitive set~\(X\) and an isomorphism \(\phi:(M,E)\cong(X,\in)\).
\def\ProblemHeadingText{Problem 5.}
\item
Sketch a proof of \(\operatorname{Con}(ZFC)\Longrightarrow \operatorname{Con}(ZFC+\neg CH+\omega_1=\omega_1^L)\). You may assume general theorems about forcing, eg. the Truth Lemma.
\def\ProblemHeadingText{Problem 6.}
\item
Define the constructible universe~\(L\) and prove that for all~\(\alpha\), \(L_\alpha\) is transitive.
\end{examproblems}
\subsection{III. Recursion Theory}
\begin{examproblems}{6}{H3}
\def\ProblemHeadingText{Problem 7.}
\item
Suppose that~\(A\) is infinite. Prove that~\(A\) is recursive if and only if there is a strictly increasing onto \(f:\mathbb{N}\to A\).
\def\ProblemHeadingText{Problem 8.}
\item
Prove that there exist recursively enumerable, recursively inseparable sets.
\def\ProblemHeadingText{Problem 9.}
\item
Prove that there is no primitive recursive function \(F(e,n)\) such that for every primitive recursive function \(f(n)\) there is an~\(e\) such that \(F(e,n)=f(n)\) for all~\(n\).
\end{examproblems}
\end{document}
